\chapter{Incremental Visualization Recommendation based on Feature Selection}

Visualization Recommendations play a crucial role in facilitating data-driven decision-making processes. These recommendations help users to select the most appropriate visualizations for a dataset. This ensures that information is presented in a clear, accurate, and precise manner, promoting better understanding and insight extraction. Effective visualization recommendations also consider the user's interest and intended applications. In an era where data is abundant, data visualization recommendations are helpful to achieve accurate results and require a better understanding of any dataset extensively to make decisions. 

According to Maruta et al. \cite{maruta2021intent}, Visualization Recommendations for the dataset are categorized into Machine Learning, Rule-based and Knowledge Graph approaches \cite{li2021kg4vis}. These approaches often struggle with scalability and complexity, particularly when handling high-dimensional data or large-scale knowledge graphs. Rule-based approaches \cite{mackinlay1986automating, mackinlay2007show, roth1994interactive, wongsuphasawat2017voyager, stolte2002polaris} have a limitation as they require expert intervention. Machine Learning approaches \cite{qian2020ml, hu2019vizml, dibia2019data2vis, luo2018deepeye, moritz2018formalizing} require dataset visualization pairs, and obtaining those pairs will be very difficult. Knowledge Graph approaches require a huge amount of data to construct the knowledge graph, which is a major hindrance. 

To address the above-discussed limitations, Maruta et al. \cite{maruta2021intent} proposed an alternate approach based on users requirements termed \textbf{Intent}. Intent is used for the dataset to filter and choose the columns from the dataset to visualize the data. This approach of using user intent to recommend visualizations can be problematic if the users intent is not specified accurately or if the dataset includes several features that correspond to the same intent, resulting in ambiguity.
   
Therefore, selecting the visualized columns to recommend the visualization for the dataset without using the users intent was quite challenging. To address these limitations, we plan to propose a Feature Selection based approach to select the columns to be visualized from the dataset, which were used to recommend the visualization without using the Intent.

Visualization Recommendation for the incremental data with respect to Volume and Velocity in terms of Big Data generated over time is critical to handle. One of our works Kuruva et al. \cite{kuruva2023visualization}, handles the incremental data exclusively and recommends the appropriate visualization. This approach adopts the Intent \cite{maruta2021intent} for existing and incremental data, and the visualization for the entire dataset is further recommended. According to \cite{kuruva2023visualization}, Intent based approach \cite{maruta2021intent} fails when the Intent is not specified by the user, which sometimes results in an incorrect recommendation. There is a demand for an alternative approach that doesn't use any Intent but still recommends accurate visualization for the Incremental data. To address these limitations, we plan to propose an Incremental Visualization Recommendation using Feature Selection \textbf{(IVR\_FS)} where the new features are considered Incremental data. 


